\subsection{A function, its domain, and its values}

\begin{corebox}[A function is more than a formula]
A function is a controlled assignment: every allowed input receives exactly one output. A formula may describe that assignment, but tables, graphs, verbal rules, and recursive rules may describe functions as well.
\end{corebox}

\begin{definitionbox}[Function]
Let $X$ and $Y$ be sets. A function
\[
f:X\to Y
\]
assigns to every $x\in X$ exactly one value $f(x)\in Y$. The set $X$ is the domain and $Y$ is the codomain.
\end{definitionbox}

The domain is part of the function. The same expression may define different functions when different inputs are allowed. For example,
\[
f(x)=\frac{1}{x-2}
\]
is meaningful only for $x\neq2$.

\begin{examplebox}[Reading values]
For
\[
f(x)=x^2+1,
\]
we have
\[
f(3)=10,
\qquad
f(-2)=5.
\]
These equations describe two values of the function, not the whole function.
\end{examplebox}

\begin{interpretationbox}[Exactly one output]
The phrase ``exactly one'' excludes two failures: assigning several outputs to one input and leaving an allowed input without an output. This precision makes expressions such as $f(a)$ unambiguous.
\end{interpretationbox}

\begin{summarybox}[Reading a function]
Identify the domain, the rule of assignment, the output corresponding to each input, and the way the function is represented.
\end{summarybox}
